% ---------------------------------------------------------------------
%  CHƯƠNG 4 — CÀI ĐẶT VÀ THỬ NGHIỆM
% ---------------------------------------------------------------------
\chapter{Cài đặt và thử nghiệm}
\label{ch:caidat}

\begin{muctieu}
Chương này mô tả môi trường, cấu trúc gói và các đoạn mã cốt lõi thể hiện kỹ
thuật bắt buộc, kèm cách chạy demo.
\end{muctieu}

\section{Môi trường triển khai}

Hệ thống phát triển trên Windows~11, JDK~21 (Temurin), chỉ dùng thư viện chuẩn.
Mã nguồn quản lý bằng \code{git} và tổ chức theo gói dưới thư mục \code{src/}.

\section{Cấu trúc gói}

\begin{lstlisting}[language=bash, caption={Cấu trúc thư mục mã nguồn.}]
src/vn/edu/dau/bxh/
  model/      Tran, Doi, DongXepHang, Lenh, YeuCau
  exception/  LoiGiaoThucException, KetQuaKhongHopLeException
  core/       BangXepHang, QuyTacXepHang
  protocol/   GiaoThuc
  luutru/     LuuTru
  server/     MayChuBXH, PhienKhach
  client/     KhachXem, TrongTai
  test/       TestXepHang, TestDongThoi
\end{lstlisting}

\section{Kỹ thuật bắt buộc và đoạn mã}

\textbf{Comparator thenComparing.} Ba tiêu chí điểm, hiệu số, bàn thắng được
ghép bằng \code{thenComparing}; cùng ba chỉ số thì lấy mã đội cho ổn định.

\begin{lstlisting}[language=Java, caption={Quy tắc xếp hạng bằng thenComparing.}]
public static final Comparator<DongXepHang> THEO_CHI_SO =
        Comparator.comparingInt(DongXepHang::diem).reversed()
                .thenComparing(Comparator.comparingInt(DongXepHang::hieuSo).reversed())
                .thenComparing(Comparator.comparingInt(DongXepHang::banThang).reversed())
                .thenComparing(DongXepHang::maDoi);
\end{lstlisting}

\textbf{Record trận và xác thực.} \code{Tran} là record bất biến; hàm tạo
compact ném ngoại lệ tự định nghĩa khi dữ liệu vô lý.

\begin{lstlisting}[language=Java, caption={Record Tran với xác thực trong compact constructor.}]
public record Tran(String doiNha, String doiKhach, int banNha, int banKhach) {
    public Tran {
        if (doiNha.equalsIgnoreCase(doiKhach))
            throw new KetQuaKhongHopLeException("Mot doi khong dau voi chinh minh");
        if (banNha < 0 || banKhach < 0)
            throw new KetQuaKhongHopLeException("Ti so khong duoc am");
    }
}
\end{lstlisting}

\textbf{Map đội và synchronized.} \code{BangXepHang} giữ \code{Map} đội và lịch
sử; mỗi thao tác ghi đều \code{synchronized}.

\begin{lstlisting}[language=Java, caption={Nhập kết quả an toàn luồng.}]
public synchronized void nhapKetQua(Tran tran) {
    dsDoi.computeIfAbsent(tran.doiNha(), ten -> new Doi(ten, ten)).ghiNhan(tran, true);
    dsDoi.computeIfAbsent(tran.doiKhach(), ten -> new Doi(ten, ten)).ghiNhan(tran, false);
    lichSu.add(tran);
}
\end{lstlisting}

\textbf{CopyOnWriteArrayList khách xem và socket đa luồng.} Máy chủ nhận kết
nối bằng \code{ServerSocket}, mỗi khách một luồng qua \code{ExecutorService};
broadcast đi qua \code{CopyOnWriteArrayList} dưới một khoá riêng.

\begin{lstlisting}[language=Java, caption={Phát broadcast an toàn cho khách xem.}]
private final CopyOnWriteArrayList<PrintWriter> dsKhachXem = new CopyOnWriteArrayList<>();
private final Object khoaGui = new Object();

void phatBangXepHang() {
    String goi = GiaoThuc.goiBangXepHang(bang.bangXepHang());
    synchronized (khoaGui) {
        for (PrintWriter ra : dsKhachXem) {
            ra.print(goi); ra.flush();
            if (ra.checkError()) dsKhachXem.remove(ra);
        }
    }
}
\end{lstlisting}

\textbf{try-with-resources khi mở tệp và socket.} \code{LuuTru} mở tệp CSV bằng
\code{try-with-resources}; \code{PhienKhach} mở socket và luồng nhập xuất cùng
cách, đảm bảo đóng ngay cả khi có ngoại lệ.

\section{Giao thức và demo}

Biên dịch rồi chạy ba tiến trình: máy chủ, một trọng tài và một khán giả.

\begin{lstlisting}[language=bash, caption={Biên dịch và chạy.}]
javac -encoding UTF-8 -d out $(find src -name '*.java')
java -cp out vn.edu.dau.bxh.server.MayChuBXH 5000 0
java -cp out vn.edu.dau.bxh.client.KhachXem 127.0.0.1 5000
java -cp out vn.edu.dau.bxh.client.TrongTai 127.0.0.1 5000
\end{lstlisting}

Trong client trọng tài gõ:
\begin{lstlisting}[language=bash]
SUBMIT HaNoi HaiPhong 2 1
SUBMIT DaNang Hue 0 0
STANDINGS
HISTORY HaNoi HaiPhong
SCHEDULE HaNoi HaiPhong DaNang Hue
EXPORT
QUIT
\end{lstlisting}

\section{Kiểm thử}

Có hai lớp kiểm thử không cần mạng. \code{TestXepHang} kiểm tra thứ tự xếp hạng,
phá đồng hạng bằng đối đầu, lịch sử đối đầu và việc chặn dữ liệu vô lý (tỉ số
âm, cùng đội, tên có khoảng trắng). \code{TestDongThoi} tạo nhiều luồng trọng
tài gửi đồng thời và kiểm tra không mất trận cùng cân đối các tổng. Chi tiết số
liệu ở Chương~\ref{ch:ketqua}.
